\section{问题重述与分析}

\subsection{问题背景与任务}

药材热风烘干是把切成圆柱形的鲜药材置于恒定或时变的热风环境中，通过表面对流同时
带走热量与水分，使内部含水率逐步降到贮藏安全线的过程。药材可近似为半径
$R_0=2\,\mathrm{cm}$、长 $L=25\,\mathrm{cm}$ 的实心圆柱；因 $L\gg R_0$，端面效应可忽略，
热质输运沿径向占主导，问题退化为一维轴对称柱坐标下的耦合传热传质过程。初始温度
$T_0=28\,^{\circ}\mathrm{C}$、初始干基含水率 $C_0=2.55$，安全含水率阈值 $C_{\mathrm{th}}=0.15$。
题目给出三套物性附录（附录2 常物性、附录3 变物性、附录4 变物性且伴随体积收缩）与两份
实测附件（附件1 热风温度与空气含湿量时序、附件2 半径随时间的收缩数据），要求依次完成：

\begin{enumerate}
\item \textbf{问题一}，在附录2 常物性、恒定热风条件下，给出预热阶段（$0\sim1800\,\mathrm{s}$）
      药材内部温度场与含水率场的时空分布；
\item \textbf{问题二}，改用附录3 温度—含水率双向耦合的变物性，刻画 $0\sim3\,\mathrm{h}$ 内
      温度场与含水率场的全过程演化；
\item \textbf{问题三}，结合附件1 的时变热风，求内部各点含水率首次全部降至 $C_{\mathrm{th}}$
      以下所需的干燥时长；
\item \textbf{问题四}，进一步计入附件2 观测到的体积收缩，建立随时间缩小的移动边界模型，
      重新给出达标时长，并与问题三比较收缩的影响。
\end{enumerate}

\subsection{整体建模思路}

四问共享同一条守恒律骨架：由能量守恒与水分质量守恒，在柱坐标下得到一对以温度 $T$ 与
干基含水率 $C$ 为状态量的非稳态抛物型方程，表面施加 Robin 型对流换热与对流传质边界，
中心施加轴对称零通量条件\upcite{dincer2010,chaedir2019}。四问的差别只体现在\textbf{物性是否随场变化}、
\textbf{边界驱动是否时变}以及\textbf{求解域是否随时间收缩}三个层次的递进，因而适合"先建公共内核、
再逐问增量扩展"的组织方式：把守恒型离散格式、Kirchhoff 界面平均与时间推进统一放入公共章
（第\ref{sec:kernel}节），各问题章只描述本问新增或改变的部分。由于表面既是升温入口又是
失水出口，轴心因对称而通量为零、响应最慢，故"内部各点均达标"的判据必然由轴心含水率控制，
这一观察贯穿后两问的终点时刻计算。以下按题号分别说明各问的难点与建模路线。

\subsection{问题一的分析}

问题一在附录2 常物性、恒定热风下刻画预热阶段的耦合场。难点有二：一是柱坐标轴心处
散度算子含 $1/r$ 奇点，需要一种在轴心天然退化、不必人工正则化的离散方式；二是传热与
传质虽同为抛物型，却在此阶段呈现截然不同的推进深度。本问作为全篇基准，重点在于给出
温度场与含水率场的时空分布，并据此判断内部是否已进入实质干燥。

\subsection{问题二的分析}

问题二改用附录3 的温度—含水率变物性，$\rho,c_p,k,D$ 均随场演化，形成双向强耦合。
难点在于非线性物性与两个守恒方程的联立求解：温度经物性反过来影响扩散，含水率又改变
热容与导热。需在同一迭代回路内联立更新 $T$ 与 $C$，并辨识传热与传质的时间尺度差异，
为后续把终点判据只建立在含水率场上提供依据。

\subsection{问题三的分析}

问题三引入附件1 的时变热风边界，求内部各点含水率首次全部降至安全线以下的时长。
难点在于时变边界驱动下的首次穿越型终点判据，以及扩散系数随含水率跨越约一个数量级
所导致的界面平均失真——简单的算术或调和平均会显著偏离真实通量。需采用守恒型的界面
处理，才能避免干燥曲线后段出现非物理的拖尾。

\subsection{问题四的分析}

问题四进一步计入附件2 观测到的体积收缩，求解域随时间缩小，成为移动边界问题。难点在于
收缩带来的动网格困难与几何度规的时变。宜借助干物质总量守恒，把移动边界的几何效应
凝聚到固定计算域上的一个时变因子中，从而避免直接在收缩物理域上离散；同时须用观测数据
独立检验"纯径向收缩"这一关键假设，以免循环论证。
